% Options for packages loaded elsewhere
% Options for packages loaded elsewhere
\PassOptionsToPackage{unicode}{hyperref}
\PassOptionsToPackage{hyphens}{url}
%
\documentclass[
  11pt,
  letterpaper,
  DIV=11,
  numbers=noendperiod,
  twoside]{scrartcl}
\usepackage{xcolor}
\usepackage[left=1.1in, right=1in, top=0.8in, bottom=0.8in,
paperheight=9.5in, paperwidth=7in, includemp=TRUE, marginparwidth=0in,
marginparsep=0in]{geometry}
\usepackage{amsmath,amssymb}
\setcounter{secnumdepth}{3}
\usepackage{iftex}
\ifPDFTeX
  \usepackage[T1]{fontenc}
  \usepackage[utf8]{inputenc}
  \usepackage{textcomp} % provide euro and other symbols
\else % if luatex or xetex
  \usepackage{unicode-math} % this also loads fontspec
  \defaultfontfeatures{Scale=MatchLowercase}
  \defaultfontfeatures[\rmfamily]{Ligatures=TeX,Scale=1}
\fi
\usepackage{lmodern}
\ifPDFTeX\else
  % xetex/luatex font selection
  \setmathfont[]{Garamond-Math}
\fi
% Use upquote if available, for straight quotes in verbatim environments
\IfFileExists{upquote.sty}{\usepackage{upquote}}{}
\IfFileExists{microtype.sty}{% use microtype if available
  \usepackage[]{microtype}
  \UseMicrotypeSet[protrusion]{basicmath} % disable protrusion for tt fonts
}{}
\usepackage{setspace}
% Make \paragraph and \subparagraph free-standing
\makeatletter
\ifx\paragraph\undefined\else
  \let\oldparagraph\paragraph
  \renewcommand{\paragraph}{
    \@ifstar
      \xxxParagraphStar
      \xxxParagraphNoStar
  }
  \newcommand{\xxxParagraphStar}[1]{\oldparagraph*{#1}\mbox{}}
  \newcommand{\xxxParagraphNoStar}[1]{\oldparagraph{#1}\mbox{}}
\fi
\ifx\subparagraph\undefined\else
  \let\oldsubparagraph\subparagraph
  \renewcommand{\subparagraph}{
    \@ifstar
      \xxxSubParagraphStar
      \xxxSubParagraphNoStar
  }
  \newcommand{\xxxSubParagraphStar}[1]{\oldsubparagraph*{#1}\mbox{}}
  \newcommand{\xxxSubParagraphNoStar}[1]{\oldsubparagraph{#1}\mbox{}}
\fi
\makeatother


\usepackage{longtable,booktabs,array}
\newcounter{none} % for unnumbered tables
\usepackage{calc} % for calculating minipage widths
% Correct order of tables after \paragraph or \subparagraph
\usepackage{etoolbox}
\makeatletter
\patchcmd\longtable{\par}{\if@noskipsec\mbox{}\fi\par}{}{}
\makeatother
% Allow footnotes in longtable head/foot
\IfFileExists{footnotehyper.sty}{\usepackage{footnotehyper}}{\usepackage{footnote}}
\makesavenoteenv{longtable}
\usepackage{graphicx}
\makeatletter
\newsavebox\pandoc@box
\newcommand*\pandocbounded[1]{% scales image to fit in text height/width
  \sbox\pandoc@box{#1}%
  \Gscale@div\@tempa{\textheight}{\dimexpr\ht\pandoc@box+\dp\pandoc@box\relax}%
  \Gscale@div\@tempb{\linewidth}{\wd\pandoc@box}%
  \ifdim\@tempb\p@<\@tempa\p@\let\@tempa\@tempb\fi% select the smaller of both
  \ifdim\@tempa\p@<\p@\scalebox{\@tempa}{\usebox\pandoc@box}%
  \else\usebox{\pandoc@box}%
  \fi%
}
% Set default figure placement to htbp
\def\fps@figure{htbp}
\makeatother


% definitions for citeproc citations
\NewDocumentCommand\citeproctext{}{}
\NewDocumentCommand\citeproc{mm}{%
  \begingroup\def\citeproctext{#2}\cite{#1}\endgroup}
\makeatletter
 % allow citations to break across lines
 \let\@cite@ofmt\@firstofone
 % avoid brackets around text for \cite:
 \def\@biblabel#1{}
 \def\@cite#1#2{{#1\if@tempswa , #2\fi}}
\makeatother
\newlength{\cslhangindent}
\setlength{\cslhangindent}{1.5em}
\newlength{\csllabelwidth}
\setlength{\csllabelwidth}{3em}
\newenvironment{CSLReferences}[2] % #1 hanging-indent, #2 entry-spacing
 {\begin{list}{}{%
  \setlength{\itemindent}{0pt}
  \setlength{\leftmargin}{0pt}
  \setlength{\parsep}{0pt}
  % turn on hanging indent if param 1 is 1
  \ifodd #1
   \setlength{\leftmargin}{\cslhangindent}
   \setlength{\itemindent}{-1\cslhangindent}
  \fi
  % set entry spacing
  \setlength{\itemsep}{#2\baselineskip}}}
 {\end{list}}
\usepackage{calc}
\newcommand{\CSLBlock}[1]{\hfill\break\parbox[t]{\linewidth}{\strut\ignorespaces#1\strut}}
\newcommand{\CSLLeftMargin}[1]{\parbox[t]{\csllabelwidth}{\strut#1\strut}}
\newcommand{\CSLRightInline}[1]{\parbox[t]{\linewidth - \csllabelwidth}{\strut#1\strut}}
\newcommand{\CSLIndent}[1]{\hspace{\cslhangindent}#1}



\setlength{\emergencystretch}{3em} % prevent overfull lines

\providecommand{\tightlist}{%
  \setlength{\itemsep}{0pt}\setlength{\parskip}{0pt}}



 


\setlength\heavyrulewidth{0ex}
\setlength\lightrulewidth{0ex}
\usepackage[automark]{scrlayer-scrpage}
\clearpairofpagestyles
\cehead{
  Brian Weatherson
  }
\cohead{
  Guessing the Less Probable
  }
\ohead{\bfseries \pagemark}
\cfoot{}
\makeatletter
\newcommand*\NoIndentAfterEnv[1]{%
  \AfterEndEnvironment{#1}{\par\@afterindentfalse\@afterheading}}
\makeatother
\NoIndentAfterEnv{itemize}
\NoIndentAfterEnv{enumerate}
\NoIndentAfterEnv{description}
\NoIndentAfterEnv{quote}
\NoIndentAfterEnv{equation}
\NoIndentAfterEnv{longtable}
\NoIndentAfterEnv{abstract}
\renewenvironment{abstract}
 {\vspace{-1.25cm}
 \quotation\small\noindent\emph{Abstract}:}
 {\endquotation}
\newfontfamily\tfont{EB Garamond}
\addtokomafont{disposition}{\rmfamily}
\addtokomafont{descriptionlabel}{\rmfamily}
\addtokomafont{title}{\normalfont\itshape}
\let\footnoterule\relax

\makeatletter
\renewcommand{\@maketitle}{%
  \newpage
  \null
  \vskip 2em%
  \begin{center}%
  \let \footnote \thanks
    {\itshape\huge\@title \par}%
    \vskip 0.5em%  % Reduced from default
    {\large
      \lineskip 0.3em%  % Reduced from default 0.5em
      \begin{tabular}[t]{c}%
        \@author
      \end{tabular}\par}%
    \vskip 0.5em%  % Reduced from default
    {\large \@date}%
  \end{center}%
  \par
  }
\makeatother
\RequirePackage{lettrine}

\renewenvironment{abstract}
 {\quotation\small\noindent\emph{Abstract}:}
 {\endquotation\vspace{-0.02cm}}

\setmainfont{EB Garamond Math}[
  BoldFont = {EB Garamond SemiBold},
  ItalicFont = {EB Garamond Italic},
  RawFeature = {+smcp},
]

\newfontfamily\scfont{EB Garamond Regular}[RawFeature=+smcp]
\renewcommand{\textsc}[1]{{\scfont #1}}

\renewcommand{\LettrineTextFont}{\scfont}
% Number theorem-like environments straight through (Theorem 1, Theorem 2, ...)
% rather than by section (1.1, 1.2, ...), which is Quarto's LaTeX default.
% Guarded, since Quarto only defines a counter for environments the document uses.
\makeatletter
\AtBeginDocument{%
  \@for\bw@thm:={theorem,lemma,corollary,proposition,conjecture,definition,%
                 example,exercise,refremark,refsolution}\do{%
    \@ifundefined{c@\bw@thm}{}{\expandafter\counterwithout\expandafter{\bw@thm}{section}}%
  }%
}
\makeatother
\KOMAoption{captions}{tableheading}
\makeatletter
\@ifpackageloaded{caption}{}{\usepackage{caption}}
\AtBeginDocument{%
\ifdefined\contentsname
  \renewcommand*\contentsname{Table of contents}
\else
  \newcommand\contentsname{Table of contents}
\fi
\ifdefined\listfigurename
  \renewcommand*\listfigurename{List of Figures}
\else
  \newcommand\listfigurename{List of Figures}
\fi
\ifdefined\listtablename
  \renewcommand*\listtablename{List of Tables}
\else
  \newcommand\listtablename{List of Tables}
\fi
\ifdefined\figurename
  \renewcommand*\figurename{Figure}
\else
  \newcommand\figurename{Figure}
\fi
\ifdefined\tablename
  \renewcommand*\tablename{Table}
\else
  \newcommand\tablename{Table}
\fi
}
\@ifpackageloaded{float}{}{\usepackage{float}}
\floatstyle{ruled}
\@ifundefined{c@chapter}{\newfloat{codelisting}{h}{lop}}{\newfloat{codelisting}{h}{lop}[chapter]}
\floatname{codelisting}{Listing}
\newcommand*\listoflistings{\listof{codelisting}{List of Listings}}
\makeatother
\makeatletter
\makeatother
\makeatletter
\@ifpackageloaded{caption}{}{\usepackage{caption}}
\@ifpackageloaded{subcaption}{}{\usepackage{subcaption}}
\makeatother
\usepackage{bookmark}
\IfFileExists{xurl.sty}{\usepackage{xurl}}{} % add URL line breaks if available
\urlstyle{same}
% fallback for those not using the hyperref driver hyperxmp:
\makeatletter
\@ifundefined{xmpquote}{\newcommand{\xmpquote}[1]{#1}}{}
\makeatother
\hypersetup{
  pdftitle={Guessing the Less Probable},
  pdfauthor={Brian Weatherson},
  hidelinks,
  pdfcreator={LaTeX via pandoc}}


\title{Guessing the Less Probable}
\author{Brian Weatherson}
\date{2026}
\begin{document}
\maketitle
\begin{abstract}
A common assumption in recent work on guessing is that it is best to
guess the most probable answer to the question asked. This note presents
three cases where this assumption looks improbable.
\end{abstract}


\setstretch{1.1}
There has been a flurry of recent interest in \emph{guessing}. Most of
this work has taken for granted something like the following principle:

\begin{description}
\tightlist
\item[Guess the Most Probable]
If one is asked \emph{Q\textsuperscript{?}}, and responds by guessing a
complete answer, that answer is appropriate iff it is the most probable
answer, given one's credences, to the question.
\end{description}

Something like \textbf{Guess the Most Probable} is a consequence of
views defended by Sophie Horowitz
(\citeproc{ref-Horowitz2019}{2019}),\footnote{Horowitz only commits to
  this for cases where there are only two possible answers, but doesn't
  give any indication that the principle won't generalise further.}
Kevin Dorst and Matthew Mandelkern
(\citeproc{ref-DorstMandelkern2022}{2022}), Ben Holguín
(\citeproc{ref-Holguin2022}{2022}), Mattias Skipper
(\citeproc{ref-Skipper2023}{2023}), Richard Teague
(\citeproc{ref-Teague2024}{2024}), and Niels Linnemann and Feraz Azhar
(\citeproc{ref-LinnemannAzhar2025}{2025}). A puzzle that Helena Fang
(\citeproc{ref-Fang2026}{2026}) raises for a generalisation of
\textbf{Guess the Most Probable} might make one suspect that there will
be problem cases for it. The point of this note is to build on Fang's
idea and describe three cases where \textbf{Guess the Most Probable}
fails. In all cases, some guesses which are not the most probable are
permitted, and indeed in all three the most probable guess has some
notable shortcomings.\footnote{This is not to say that it is
  impermissible; the cases are hard, and any answer might seem
  defective. But in all cases the most probable answer seems weird
  enough that it permits alternative answers.}

\section{Guess the Order}\label{sec-order}

At the start of every Premier League season, the \emph{Guardian Football
Weekly} asks each of its panellists to guess the order that the 20 teams
will finish in. This is obviously absurd; there are 20! possible orders,
and while some can be ruled out fairly comprehensively, there are too
many reasonable ones to guess any with any confidence. But guess they
do, and it is a useful framing device for talking about the upcoming
season.

Inspired by that, I started looking at guessing the order of finishers
in somewhat smaller competitions. So consider a simple four-team league,
where each team plays each other once, not twice as in the Premier
League. You're probably familiar with these competitions; they are the
standard form of the group stage at major international tournaments. We
have four teams, imaginatively called A, B, C and D, and our expert is
asked to guess the order that they'll finish in.

Our expert has to deal with two kinds of uncertainty.

First, she doesn't know exactly how good each team is. She knows each
team has an underlying quality, and for each team she has a distribution
over their possible quality values, but she could be wrong. We'll
simplify greatly; in each case the probability distribution is a normal
distribution with identifiable mean and variance. Crucially, she is very
unsure about one of the teams, namely C. For that team, her credence
about their quality has a very high variance, while she is fairly
confident about the other three.

Second, soccer games are just random. You may remember Saudi Arabia
beating Argentina in the 2022 World Cup. That was not because Saudi
Arabia was the better team; sometimes the better team loses. The closer
the teams are in quality, the more likely an upset is. But even for
quite large quality differences, like Saudi Arabia and Argentina, the
upset is possible.

Here's the model this paper uses for this kind of guessing.\footnote{AI
  acknowledgment: I used Claude Opus 5.5 to build this model, having
  explained to it the kind of tournament that I wanted. Once it was
  done, it came up unprompted with some summary statistics about the
  model which in turn I use below.}

Each team \(i\) has a real quality \(q_i\). In each tournament
simulation, the qualities are drawn from the normal distributions with
means and standard deviations described in Table~\ref{tbl-quality}. In
any game, we start with the difference \(d\) between the two team
qualities (remember quality is held fixed across a tournament). If the
home team\footnote{`Home team' here is just the first listed team. We
  are not assuming any home field advantage.} is better, \(d\) is
positive; if they are worse \(d\) is negative. Then we add a noise
parameter drawn from \(\mathcal{N}(0, 1)\) and tripled. If the absolute
value of the result is less than 1.2, the game is a draw. Otherwise, if
the result is positive the home team wins, and if it is negative the
away team wins.

{

\begin{longtable}[]{@{}lrr@{}}

\caption{\label{tbl-quality}The expert's credences about team quality}

\tabularnewline

\toprule\noalign{}
Team & Mean & Standard deviation \\
\midrule\noalign{}
\endhead
\bottomrule\noalign{}
\endlastfoot
A & 4.00 & 0.2 \\
B & 3.00 & 0.2 \\
C & 2.00 & 1.6 \\
D & 1.55 & 0.2 \\

\end{longtable}

}

To see what the probabilities were for various outcomes, we then
simulated the season two million times. As in regular soccer, a team got
three points for a win, and one point for a draw. After three games, we
used a randomiser to split teams on equal points.\footnote{We could have
  tried to simulate not just results but scores, then used tie-breaking
  procedures, but this was getting complicated enough.} And we measured
how often each of the 24 possible permutations of the teams came up over
those runs.

Now this model obviously leaves out a lot. It doesn't allow for team
quality changing over time, due to injuries. It doesn't allow for teams
matching up particularly well or badly against other teams. But it's
complicated enough to serve as an illustration.

Assume our expert knows her own mind, so she knows the numbers in
Table~\ref{tbl-quality}. And she knows how team quality relates to win
probabilities, via noise parameters. She has just as many computational
resources as we do, so she can build all the tables you'll see here. And
she is asked, just like the podcast panellists, to guess the order in
which teams will finish in this group.

If she gives a complete answer, and really that's what panellists are
supposed to do, and \textbf{Guess the Most Probable} is correct, then
there is only one permissible answer. It's ABDC. The numbers are shown
in Table~\ref{tbl-rank-frequency}.

{

\begin{longtable}[]{@{}rlr@{}}

\caption{\label{tbl-rank-frequency}The six most probable finishing
orders}

\tabularnewline

\toprule\noalign{}
Rank & Order & Probability \\
\midrule\noalign{}
\endhead
\bottomrule\noalign{}
\endlastfoot
1 & ABDC & 0.136 \\
2 & ABCD & 0.128 \\
3 & ACBD & 0.093 \\
4 & BADC & 0.082 \\
5 & BACD & 0.079 \\
6 & ADBC & 0.072 \\

\end{longtable}

}

But that's a surprising guess for at least three reasons.

First, she thinks that C is probably better than D. The mean of C's
quality distribution is higher than the mean of D's. It's completely
reasonable to rank teams in a task like this by your mean estimate of
their quality. (I suspect that's more or less what people really do when
asked to produce these rankings.)

Second, our expert will do better at predicting the relative finishing
positions of pairs of teams if she guesses ABCD than if she guesses
ABDC. There are six pairs of teams in the tournament, and we can ask for
each pair who finished higher in any tournament run. For each pair, we
can also ask whether the actual order of those two teams matches the
order the expert predicted for them. It turns out, as shown in
Table~\ref{tbl-pair-match}, that in expectation she matches more pairs
by guessing ABCD than guessing ABDC. Indeed, ABCD does better on this
measure than any other order.

{

\begin{longtable}[]{@{}lrrr@{}}

\caption{\label{tbl-pair-match}Expected number of pairs in the right
order, and of teams in the right place, for the six most probable
finishing orders}

\tabularnewline

\toprule\noalign{}
Order & Probability & Expected pairs & Expected places \\
\midrule\noalign{}
\endhead
\bottomrule\noalign{}
\endlastfoot
ABDC & 0.136 & 4.054 & 1.537 \\
ABCD & 0.128 & 4.182 & 1.563 \\
ACBD & 0.093 & 3.913 & 1.420 \\
BADC & 0.082 & 3.737 & 1.224 \\
BACD & 0.079 & 3.866 & 1.249 \\
ADBC & 0.072 & 3.606 & 1.305 \\

\end{longtable}

}

The point here is not that thinking about relative finishing positions
of pairs raises extra questions over and above the question of what the
finishing order will be. To guess ABCD is, inter alia, to guess that C
will finish above D. All that we need to make ABCD a sensible guess is
that the guesser cares about these extra things that she commits herself
to when she guesses an order.

There is a natural way to make this verdict fit with the picture that
Dorst and Mandelkern give. At one point in their paper, they are
interested in how the guesser might trade off between more and less
specific guesses. (For example, when asked to guess who will win the
group, when is it best to say A, and when is it best to say something
less committal like A or B?) They use a standard epistemic utility model
to model these decisions, where one of the factors is how accurate the
guess is. Let \(v_A(g, w)\) be the accuracy value of guess \(g\) in
world \(w\). Dorst and Mandelkern make the natural assumption that
\(v_A(g, w) = 1\) if \(g\) is true in \(w\), and 0 otherwise. They also
note that you might want a more complicated assumption if you care about
verisimilitude or partial truth.

That provides a natural explanation for why ABCD may be a permissible
guess even though ABDC is more probable. Say \(v_A(g, w)\) is the number
of pairs whose finishing order is the same in \(g\) and \(w\). (So if
\(g\) is completely correct in \(w\), \(v_A(g, w) = 6\); if it just gets
one pair flipped around, it's 5, and so on.) Then the guess with the
highest expected accuracy, by this accuracy measure, will be ABCD. The
expected pairs column in Table~\ref{tbl-pair-match} just is the expected
accuracy score by this measure. In this case the guess which orders the
teams by their mean quality will be a Condorcet winner. That is, one
team will be guessed to finish higher than another iff the probability
that it finishes higher is above 0.5.

Third, although it is very close, we can rerun the second argument if
the guesser cares not about the relative positions of teams, but their
absolute finishing position. Another natural setting for \(v_A(g, w)\)
sets it equal to how many teams finish in the same position in \(g\) and
\(w\). That's the right metric for a guesser who really wants to
correctly guess who will finish first, and who will finish second, and
so on. If we do that, the expected places column in
Table~\ref{tbl-pair-match} shows that ABCD has, by a whisker, the
highest expected accuracy. Like with the pairs, the absolute finishing
positions of the teams isn't another thing to guess above their order;
guessing ABCD just is a way to guess that C will finish third.

So there are several reasons to guess ABCD. There is still one reason to
guess ABDC; it is the most probable. And I don't deny that it is, for
that reason, a permissible guess. What I do deny is that it is the only
permissible complete guess. It is permissible to rank teams by their
mean quality, or to prefer the Condorcet ordering, or to maximise the
expected number of finishing positions one correctly predicts, even if
that means making a guess with a slightly lower probability of being
perfectly correct. As long as ABCD is one of the permissible guesses,
that is enough to refute \textbf{Guess the Most Probable}, which says
that only ABDC is.

In the case I've described so far ABCD is the second-most probable
guess; I couldn't find a model where it was best on both of the measures
listed in Table~\ref{tbl-pair-match} and was lower than second-most
probable. You might wonder whether this suggests that \textbf{Guess the
Most Probable} is at least close to true. If the argument about pairs is
correct, and it is always permissible to guess the Condorcet ranking,
then it is permissible to guess really improbable orderings. Imagine our
expert was asked to provide a best guess about the finishing order of
the group with teams E, F, G and H, where the means and standard
deviations of her credences about their quality are given in table
Table~\ref{tbl-high-variance}.

{

\begin{longtable}[]{@{}lrr@{}}

\caption{\label{tbl-high-variance}The expert's credences about team
quality for the second group}

\tabularnewline

\toprule\noalign{}
Team & Mean & Standard deviation \\
\midrule\noalign{}
\endhead
\bottomrule\noalign{}
\endlastfoot
E & 1.4 & 0.2 \\
F & 1.3 & 2.0 \\
G & 1.2 & 2.0 \\
H & 1.1 & 0.2 \\

\end{longtable}

}

As Table~\ref{tbl-high-variance-orders} shows, guessing that they will
finish EFGH maximises the expected number of pairs one expects to order
correctly. It's also the Condorcet order: whenever \(x\) is more than
50\% likely to finish about \(y\), it is guessed to finish above \(y\).
But it is not the most probable order. In fact, it is the twenty-first
most probable order out of twenty-four. The second-best guess by
expected pairs guessed correctly is the twenty-second most probable
order. Once we start caring about claims embedded within a guess, the
best guess might be a long way down the probability ranking.

{

\begin{longtable}[]{@{}lrrrr@{}}

\caption{\label{tbl-high-variance-orders}Statistics for the second
group, showing the orders that finish in the top three on some
interesting statistic or other.}

\tabularnewline

\toprule\noalign{}
Order & Rank & Probability & Expected pairs & Expected places \\
\midrule\noalign{}
\endhead
\bottomrule\noalign{}
\endlastfoot
FEHG & 1 & 0.066 & 3.085 & 1.180 \\
GEHF & 2 & 0.062 & 3.013 & 1.150 \\
FHEG & 3 & 0.059 & 2.987 & 1.134 \\
FEGH & 7 & 0.046 & 3.112 & 1.046 \\
EFHG & 8 & 0.045 & 3.111 & 1.045 \\
GEFH & 11 & 0.044 & 3.064 & 1.030 \\
EFGH & 21 & 0.030 & 3.138 & 0.911 \\
EGFH & 22 & 0.029 & 3.115 & 0.906 \\

\end{longtable}

}

That said, this second group is hard, and there are good reasons to not
guess EFGH. It does not maximise the probability that one will guess the
group winner correctly. It does maximise the probability that one will
guess who the top two teams are correctly, which is often important in
tournaments, but in that respect it's tied with more probable guesses.
While it often correctly guesses the identity of the top two, it is much
more rare that it guesses their order, i.e., who is first and who is
second, which is also important for tournaments. Just like it is the
twenty-first most likely order, it ranks twenty-first, out of twenty
four, in the number of teams it is expected to correctly guess the
position of. I still think that all things considered, EFGH is a
reasonable guess, but there are many reasons for guessing other orders.

The general claim of this section, of this paper really, is that when
asked to make a \emph{structured} guess, it is generally reasonable to
have some kind of more complicated value function than just 1 for true
and 0 for false. What's distinctive about structured guesses, like
guessing which order the teams will finish in, is that they commit the
guesser to answers to related questions. The guesser might, qua guesser,
care about these extra answers being correct. And if they do, they have
a reason purely within accuracy considerations to guess ABCD.

You might worry this will overgeneralise. Why does this have anything to
do with structure? Imagine someone asked to guess where James Bond is,
and they guess Vienna. Here is a parallel to the earlier arguments about
commitment. By guessing Vienna, they commit to Bond being within 50
miles of Vienna, so their guess should get some accuracy points if he is
in Bratislava. And they commit to Bond being in Austria, so they should
get some accuracy points if he is in Salzburg. So even for unstructured
questions like \emph{Where is James Bond?}, there could be a reason to
use something other than a 0/1 accuracy score, and hence reason to make
some guess other than the most probable.

To respond to this worry, we need a tighter notion of commitment. It's
true that Bond being in Bratislava or in Salzburg is, in different ways,
like his being in Vienna. But the inferences from \emph{Bond is in
Vienna} to \emph{Bond is near Vienna} or \emph{Bond is in Austria}
involve a change of subject-matter. The core idea I'm using here is
taken from Yablo (\citeproc{ref-Yablo2014}{2014}); see Fine
(\citeproc{ref-Fine2016}{2016}) and Berto
(\citeproc{ref-Berto2022}{2022}) for related work. The inference from a
conjunction to a conjunct, or from a ranking to a pair-ordering, does
not involve changing the subject. If you guess a conjunction, the
conjuncts are quite literally parts of the guess. So roughly (and see
the earlier references for a less rough version), the idea is that by
guessing \(g\), you are committed to guessing anything that \(g\)
entails whose subject-matter is part of \(g\)'s. That's why the
structured guesses allow for guessing the less probable; the relevant
accuracy measure won't be the simple 0/1 measure.

Finally, a housekeeping note. Dorst and Mandelkern
(\citeproc{ref-DorstMandelkern2022}{2022}, fn. 15) list two reasons for
thinking that the accuracy value might be something other than 0 or 1:
verisimilitude or partial truth. The argument of the last few paragraphs
is that partial truth really is important, but verisimilitude, i.e.,
closeness to the truth, is not. Guessing Bond is in Vienna when he's in
Bratislava is close; it's less than 50 miles away! But for the
\emph{Where is James Bond?} question, I've endorsed the claim that
guesses should just be judged true or false, and the right (complete)
guess is the most probable one. The examples that I do think make it
reasonable for the guesser to have a more complicated value function are
what they are calling partial truth cases. Not coincidentally, they cite
Yablo as a place to look for more detail on partial truth, so the kinds
of cases I'm using are pretty close to what they were considering.

\section{Follow Up Guesses}\label{sec-follow}

The next case is largely a variant on a case developed by Richard Teague
(\citeproc{ref-Teague2024}{2024}); what I want to stress is the problem
it raises for \textbf{Guess the Most Probable}.

Having settled the ABCD and EFGH groups, the expert is now thinking
about the JKLM group. The means and standard deviations of her credences
about team qualities in that group are given in
Table~\ref{tbl-quality-sec-2}. Note that the variance is now really
high; she is very unsure about all these teams.

{

\begin{longtable}[]{@{}lrr@{}}

\caption{\label{tbl-quality-sec-2}The expert's credences about team
quality in the JKLM group}

\tabularnewline

\toprule\noalign{}
Team & Mean & Standard deviation \\
\midrule\noalign{}
\endhead
\bottomrule\noalign{}
\endlastfoot
J & 4.0 & 1 \\
K & 2.5 & 4 \\
L & 1.1 & 2 \\
M & 1.0 & 3 \\

\end{longtable}

}

On a podcast previewing the tournament, they get to this group, and the
following exchange occurs:

\begin{quote}
Host: Who will win this group?\\
Expert: J. These teams are hard to judge, but J is likely best.\\
Host: Makes sense. Who will finish second?\\
Expert: {[}panics{]}
\end{quote}

The reason that the expert panics at this point is that it is really
unclear how she is supposed to answer. To see this, look at her
credences for which team will finish in which position in
Table~\ref{tbl-place-prob-sec-2}.

{

\begin{longtable}[]{@{}lrrrr@{}}

\caption{\label{tbl-place-prob-sec-2}The expert's credences about where
each team will finish}

\tabularnewline

\toprule\noalign{}
Team & First & Second & Third & Fourth \\
\midrule\noalign{}
\endhead
\bottomrule\noalign{}
\endlastfoot
J & 0.443 & 0.355 & 0.158 & 0.043 \\
K & 0.313 & 0.219 & 0.214 & 0.254 \\
L & 0.102 & 0.223 & 0.345 & 0.331 \\
M & 0.142 & 0.203 & 0.283 & 0.372 \\

\end{longtable}

}

If \textbf{Guess the Most Probable} is correct, then she should answer
J. That's the team with the highest probability of finishing second. But
that would be a real violation of the conversational norms in this
situation. Or, at least, I think it would be a pretty big breach. So
\textbf{Guess the Most Probable} is wrong.

To be sure, if this was the only problem with \textbf{Guess the Most
Probable}, then it would be fairly easy to fix it to capture what the
authors listed in the introduction want. The claim to fall back on is
that the most probable option should be guessed when one is giving a
complete answer to a question that is discourse-initial. It's arguably
the earlier guess that's the problem here.

If the fix is so easy, why did I include the case? Three reasons. First,
even if the fix is easy, it's worth making. Second, although intuitions
may vary here, I think there's something a bit funny about answering
\emph{J} if the first question is \emph{Who will finish second?}. I'm
not sure about that intuition, and I don't have a great story about how
to capture it, but it seems like an interesting puzzle. Third, the more
interesting (to me) puzzle is about what the expert should say in the
actual conversation. There are a few moves available.

She could just say the most probable team to finish second, i.e., J.
That sounds terrible to me, but given the complications involved in the
other answers, maybe that's the right thing to say.

She could say the team that's most likely to finish second consistent
with her earlier answers. That, as Table~\ref{tbl-place-prob-sec-2}
shows, is L. In some ways that's the simplest fix to \textbf{Guess the
Most Probable}; say that in guesses that appear some way into a
discourse, there is a requirement of consistency with past guesses.

This answer is not obviously optimal for similar reasons to why guessing
ABDC was not obviously optimal. For one thing, K has the higher mean
quality. For another, K has the highest probability, of the teams not
yet listed, of finishing in the top two. For yet another, and this I
think is most crucial, K has the highest probability of finishing second
conditional on the previous guess being correct. The probabilities of
each team finishing second given the earlier guess are given in
Table~\ref{tbl-conditional-second}, and K is clearly highest.

{

\begin{longtable}[]{@{}lr@{}}

\caption{\label{tbl-conditional-second}Probability of finishing second,
given that J wins the group}

\tabularnewline

\toprule\noalign{}
Team & Probability \\
\midrule\noalign{}
\endhead
\bottomrule\noalign{}
\endlastfoot
K & 0.373 \\
L & 0.332 \\
M & 0.295 \\

\end{longtable}

}

So the third option is to guess K. One motivation for this is the idea
that when asked to make a series of guesses, one is being asked to
develop a coherent story. One should maximise not the probability of the
next step, but the probability of the whole story given the possible
continuations of it. That means guessing the option with the highest
conditional probability, or, if it is a structured guess, the highest
expected accuracy.\footnote{Fang (\citeproc{ref-Fang2026}{2026}, fn. 36)
  briefly discusses this idea, and dismisses it as changing the
  question. I am somewhat more optimistic that it can work.} That seems
like a plausible story of what we're doing with sequential guesses, and
it gets what I think is the intuitive answer here.

There is a fourth option, which isn't quite inconsistent with the
previous ones. If sequential guessing is a matter of developing a story
piece-by-piece, then guessers might want the option to eventually revise
their previous guesses. After a while, a sequence of guesses might lead
to too implausible a story. The point is not just that the story you end
up with is absolutely implausible; sometimes the best guess will be very
improbable. (Think of the pundits asked to choose which of 20! orderings
of Premier League teams will be the right one.) Rather, the point is
that sometimes when you see the full set of questions you are being
asked, you can make a much more plausible set of guesses by retracting
some earlier guesses.

It doesn't matter for this example, but one thing in general that the
expert could do at the panic stage is to retract their earlier guess.
That isn't really consistent with the norms of panel discussions in mass
media, but might be consistent with the norms of guessing in less loaded
contexts. They could say ``Oh, you want me to guess the top two in that
order; in that case I'd say\ldots{}'', or ``Oh, you're going to ask for
a full ranking; in that case I'd say\ldots{}''. The view that guesses
should maximise conditional accuracy value, conditional on the previous
guesses, leads to some very implausible outcomes if you don't add in a
retraction option, so this fourth option might be best thought of as a
complement to the third, rather than an alternative.

I don't have a good answer here to what the expert should do; apart from
noting that the case suggests that \textbf{Guess the Most Probable}
needs to be at least amended, I'm mostly presenting it as a puzzle for
further research.

\section{Guesses with Incoming Information}\label{sec-incoming}

The last case is basically an application of some points Sarah Moss
(\citeproc{ref-Moss2015sp}{2015}) made about disjunctions and
context-sensitivity. Let's leave soccer and return to a case beloved in
the guessing literature: horse racing.

There is a big race tomorrow with thirteen runners. Three of them, N, O
and P, are very good; the other ten are not. The expert has views on two
questions: how well each horse runs in different conditions, and what
the conditions are likely to be. The latter is important. There is a big
storm moving towards the racetrack. It will reach us overnight. While
the storm is large, it is concentrated, so it's possible it will pass
just north of here, but it's also possible it will hit us hard. If it
passes by, the track will be dry and fast; if it hits, the track will be
a slow muddy mess. Horse N is good in all conditions, but O and P have
preferred conditions. O is a flat-track bully, able to sprint away on
dry, fast tracks. P is a mudlark, able to churn through mud better than
most. So the expert's credences about the race results are highly
dependent on what happens overnight with the storm. In
Table~\ref{tbl-dry} and Table~\ref{tbl-wet}, I've listed her credences
as to who will win conditional on a dry track, and conditional on a wet
track.

\begin{table}

\caption{\label{tbl-going}The expert's win probabilities, conditional on
the state of the track}

\begin{minipage}[t]{0.50\linewidth}

\subcaption{\label{tbl-dry}Dry track}

\centering{

\begin{tabular}{lr}
\toprule
Horse & Win probability\\
\midrule
N & 0.35\\
O & 0.45\\
P & 0.10\\
Others, each & 0.01\\
\bottomrule
\end{tabular}

}

\end{minipage}%
%
\begin{minipage}[t]{0.50\linewidth}

\subcaption{\label{tbl-wet}Wet track}

\centering{

\begin{tabular}{lr}
\toprule
Horse & Win probability\\
\midrule
N & 0.35\\
O & 0.10\\
P & 0.45\\
Others, each & 0.01\\
\bottomrule
\end{tabular}

}

\end{minipage}%

\end{table}%

The expert thinks it is 50/50 which kind of track we get; the weather
forecasts are very uncertain about the trajectory of the storm. So the
highest win probability belongs to N, at 0.35. But guessing that N will
win has some downsides. Notably, it violates a reflection principle for
guessing. If the expert guesses N, she knows that whatever happens
overnight, in the morning she'll want to change that guess. This feels
like something we should avoid, if possible. And the expert has an easy
way out. Instead of guessing N, she can say (1).

\begin{enumerate}
\def\labelenumi{(\arabic{enumi})}
\tightlist
\item
  If it is dry, O will win, or if it is wet, P will win.
\end{enumerate}

She could rephrase this as (2).

\begin{enumerate}
\def\labelenumi{(\arabic{enumi})}
\setcounter{enumi}{1}
\tightlist
\item
  It will be either dry or wet, so either O or P will win.
\end{enumerate}

Finally, and this is where I'm borrowing most heavily from Moss, she
might be in a situation where it is common ground that there is a
substantial probability both that the storm will hit and that it will
pass us by, that this fact is relevant to the win probabilities, and
that O is a flat-track specialist while P is a mudlark. Then it seems
she can say (3) as a way of expressing what she would say more
explicitly as (2) or (1).

\begin{enumerate}
\def\labelenumi{(\arabic{enumi})}
\setcounter{enumi}{2}
\tightlist
\item
  Either O or P will win.
\end{enumerate}

If, in that context, (3) is a reasonable guess, then we have a
counterexample to what Dorst and Mandelkern call \textbf{Filtering}, and
what Holguín calls \textbf{cogency}. This is a disjunctive guess that's
permissible, despite there being another guess you could generate by
swapping one of its disjuncts for a more probable complete answer,
namely N.

Of course, it isn't a counterexample to \textbf{Guess the Most Probable}
because it isn't a complete guess. But we can fix that. Let's say a new
forecast comes in and it's now 55\% likely to be a wet track. Again, the
expert is asked to predict the race, and they say (4).

\begin{enumerate}
\def\labelenumi{(\arabic{enumi})}
\setcounter{enumi}{3}
\tightlist
\item
  I think it will be a wet track, so P.
\end{enumerate}

That all seems reasonable to me. It's a less probable guess than N, but
it spells out the reason, and it is not known to be a reflection
violation. In a context where it is obvious that you would only guess P
if you thought it would be a wet track, and the wet track is 55\%
likely, I think you can even get (5) to come out appropriate.

\begin{enumerate}
\def\labelenumi{(\arabic{enumi})}
\setcounter{enumi}{4}
\tightlist
\item
  I think P; he's the best in the mud.
\end{enumerate}

And that is a violation of \textbf{Guess the Most Probable}. Now I don't
want to lean too heavily on the intuitions here; I don't think it's
nearly as clear that \emph{P} is a good guess here as that \emph{ABCD}
is a good guess to the ordering question in Section~\ref{sec-order}. All
I do want to say is that (a) it isn't obviously a bad guess, and (b) the
fact that \textbf{Guess the Most Probable} leads to guesses that one
knows one will want to retract should count against it.

\section{Conclusion}\label{conclusion}

When presented with a question, and asked to guess, guessers can
reasonably have many more goals than saying the thing that most likely
is the true answer to that question. They might try to truly answer as
many parts of the question as they can; that was the consideration
behind the main example in Section~\ref{sec-order}. They might try to
craft a coherent story given their prior guesses; that was the
consideration behind the main example in Section~\ref{sec-follow}. They
might try to make guesses that they have some chance of standing behind
when new information comes in; that was the consideration behind the
main example in Section~\ref{sec-incoming}. By looking at cases like
these, we get a much richer picture of what people are trying to do with
guesses.

\subsection*{References}\label{references}
\addcontentsline{toc}{subsection}{References}

\protect\phantomsection\label{refs}
\begin{CSLReferences}{1}{0}
\bibitem[\citeproctext]{ref-Berto2022}
Berto, Francesco. 2022. \emph{Topics of Thought: The Logic of Knowledge,
Belief, Imagination}. Oxford: Oxford University Press.

\bibitem[\citeproctext]{ref-DorstMandelkern2022}
Dorst, Kevin, and Matthew Mandelkern. 2022. {``Good Guesses.''}
\emph{Philosophy and Phenomenological Research} 105 (3): 581--618. doi:
\href{https://doi.org/10.1111/phpr.12831}{10.1111/phpr.12831}.

\bibitem[\citeproctext]{ref-Fang2026}
Fang, Helena. 2026. {``Guessing and Its Limits.''} \emph{Philosophy and
Phenomenological Research} 112 (3): 563--77. doi:
\href{https://doi.org/10.1111/phpr.70101}{10.1111/phpr.70101}.

\bibitem[\citeproctext]{ref-Fine2016}
Fine, Kit. 2016. {``Angellic Content.''} \emph{Journal of Philosophical
Logic} 45 (2): 199--226. doi:
\href{https://doi.org/10.1007/s10992-015-9371-9}{10.1007/s10992-015-9371-9}.

\bibitem[\citeproctext]{ref-Holguin2022}
Holguín, Ben. 2022. {``Thinking, Guessing, and Believing.''}
\emph{Philosophers' Imprint} 22 (6). doi:
\href{https://doi.org/10.3998/phimp.2123}{10.3998/phimp.2123}.

\bibitem[\citeproctext]{ref-Horowitz2019}
Horowitz, Sophie. 2019. {``Accuracy and Educated Guesses.''} In
\emph{Oxford Studies in Epistemology}, edited by Tamar Szabó Gendler and
John Hawthorne. Vol. 6. Oxford University Press.

\bibitem[\citeproctext]{ref-LinnemannAzhar2025}
Linnemann, Niels, and Feraz Azhar. 2025. {``Better Guesses.''}
\emph{Philosophy and Phenomenological Research} 110 (2): 661--86. doi:
\href{https://doi.org/10.1111/phpr.13118}{10.1111/phpr.13118}.

\bibitem[\citeproctext]{ref-Moss2015sp}
Moss, Sarah. 2015. {``On the Semantics and Pragmatics of Epistemic
Vocabulary.''} \emph{Semantics and Pragmatics} 8 (5): 1--81. doi:
\href{https://doi.org/10.3765/sp.8.5}{10.3765/sp.8.5}.

\bibitem[\citeproctext]{ref-Skipper2023}
Skipper, Mattias. 2023. {``Good Guesses as Accuracy-Specificity
Tradeoffs.''} \emph{Philosophical Studies} 180 (7): 2025--50. doi:
\href{https://doi.org/10.1007/s11098-022-01915-1}{10.1007/s11098-022-01915-1}.

\bibitem[\citeproctext]{ref-Teague2024}
Teague, Richard. 2024. {``A Puzzle about Guessing and Inquiry.''}
\emph{Analysis} 84 (2): 327--36. doi:
\href{https://doi.org/10.1093/analys/anad057}{10.1093/analys/anad057}.

\bibitem[\citeproctext]{ref-Yablo2014}
Yablo, Stephen. 2014. \emph{Aboutness}. Princeton: Princeton University
Press.

\end{CSLReferences}




\end{document}
